\chapter{Introduction}

Universities are large, complex organisations that provide a wide range of
services to their students. At King's College London (KCL), students regularly
need support across areas such as academic assessment, welfare, financial
matters, careers, and university procedures. While specialist services exist for
each of these areas, students frequently direct their queries to generalist
staff---such as programme officers and personal tutors---who must then triage,
redirect, or respond to a significant volume of incoming messages. This creates
a recurring burden: staff spend considerable time handling queries that are
routine, misdirected, or duplicated across multiple recipients, while overdue
and unanswered queries risk going unnoticed altogether. The absence of a
structured query-management process means that both staff efficiency and student
experience suffer.

The KCL Ticketing System addresses this problem by providing programme officers
with a centralised web application for managing incoming student queries as
support tickets. Programme officers gain a single, structured workspace where
every query is visible, trackable, and attributable---eliminating the risk of
queries going unanswered and removing the duplicated effort that arises when the
same query reaches multiple members of staff. Students benefit from a reliable,
transparent channel for raising issues, with confidence that their query has
been received and will be handled. The system further reduces staff workload by
enabling immediate, automated preliminary responses to routine queries,
reserving human attention for the cases that genuinely require it.

The system is a full-stack web application accessible through any modern
browser. The back end is built with \textbf{Django 5} and
\textbf{Django REST Framework} (Python), exposing a JSON API secured with
JWT authentication. The front end is a single-page application built with
\textbf{React 18} and \textbf{Redux Toolkit} (JavaScript). An AI-powered
support assistant uses the \textbf{Google Gemini API} to provide automated
responses to routine student queries. Data is persisted in
\textbf{SQLite} for development, with the architecture supporting migration to a
production-grade relational database.